% Generated draft: requires mathematical and source review.
\documentclass[11pt]{article}
\usepackage[margin=25mm]{geometry}
\usepackage{fontspec}
\setmainfont{Latin Modern Roman}
\usepackage{amsmath,amssymb,mathtools,unicode-math}
\setmathfont{Latin Modern Math}
\usepackage{enumitem,needspace,xurl,hyperref}
\setlength{\parindent}{0pt}
\setlength{\parskip}{5pt}
\setlength{\emergencystretch}{3em}
\newcommand{\N}{\mathbb{N}}
\newcommand{\Z}{\mathbb{Z}}
\newcommand{\Q}{\mathbb{Q}}
\newcommand{\R}{\mathbb{R}}
\newcommand{\C}{\mathbb{C}}
\newcommand{\F}{\mathbb{F}}
\newcommand{\A}{\mathbb{A}}
\newcommand{\PP}{\mathbb{P}}
\newcommand{\E}{\mathbb{E}}
\newcommand{\eps}{\varepsilon}
\newcommand{\dd}{\,\mathrm{d}}
\begin{document}
\section*{UnsolvedMath ID: 30006246}
\textbf{Source problem code:} catalogue entry
\subsection{Definitions and mathematical statement}
Let \(G\) be a finite group and let \(X\) be a specified nonempty finite set of formal variables. Write \(F(X)\) for the free group on \(X\), and \(G*F(X)\) for its free product with \(G\). For each assignment \(a:X\to G\), let \(\operatorname{ev}_a:G*F(X)\to G\) be the homomorphism that fixes every element of \(G\) and sends \(x\) to \(a(x)\). A mixed identity over \(G\) is an element \(w\in G*F(X)\) such that
\[
w\ne 1\quad\text{in }G*F(X),\qquad
\operatorname{ev}_a(w)=1_G\quad\text{for every }a:X\to G.
\]
Let \(\pi:G*F(X)\to F(X)\) be the homomorphism that sends every element of \(G\) to the identity and fixes the variables. Under the conventional interpretation used here, a mixed identity is nonsingular when \(\pi(w)\ne 1\). The supplied excerpt does not explicitly define nonsingularity, so this interpretation requires confirmation against the report.

For a specified word-length convention \(\ell\), define
\[
\begin{aligned}
m_{\mathrm{mix}}(G;X,\ell)
&=\min\{\ell(w):w\ne1,\ \operatorname{ev}_a(w)=1_G\text{ for every }a:X\to G\},\\
m_{\mathrm{ns}}(G;X,\ell)
&=\min\{\ell(w):w\ne1,\ \pi(w)\ne1,\ \operatorname{ev}_a(w)=1_G\text{ for every }a:X\to G\}.
\end{aligned}
\]
Both minima are taken over words in \(G*F(X)\); an empty minimum is assigned \(+\infty\). These expressions make the dependencies explicit without selecting an unsupported length convention or number of variables.

The supplied problem asks for the lengths of shortest mixed identities and shortest nonsingular identities for the \emph{remaining families of finite groups of Lie type}. Once the intended families, allowed variables, and length convention are specified, the requested quantities are the corresponding minima above, with minimization over variable sets as well if that is the intended convention. The supplied record does not identify those remaining families, distinguish simple groups from other forms of groups of Lie type, specify admissible rank and field parameters or exceptional small cases, define length, or say whether exact values or asymptotic orders are sought. Consequently, it does not determine a unique fully specified mathematical problem; these qualifications must remain unresolved in this statement-only draft.
\subsection{Short English statement}
Determine how short a nontrivial word with constants from a finite group of Lie type can be while evaluating to the identity for every substitution, and determine the corresponding shortest length when the word remains nontrivial after its constants are erased. The intended remaining families, variable convention, length convention, and meaning of the requested determination require confirmation from the source context. The supplied partially solved label is dated curation, not an independently verified conclusion.
\subsection{Sources}
\begin{itemize}
\item Local UnsolvedMath registry, \texttt{data/unsolvedmath/problems.json}, supplied record with permanent ID \(30006246\), titled \emph{Mixed and Nonsingular Identities in Finite Lie-Type Groups}; exact source problem code \texttt{catalogue entry}. The record attributes the problem to the 2025 Oberwolfach Reports collection and names no individual proposer.
\item Supplied report citation: \emph{Finite Groups, Fusion Systems and Applications} (2025), Oberwolfach Reports, pp.\ 689--752, DOI \texttt{10.4171/owr/2025/16}. \url{https://doi.org/10.4171/owr/2025/16}. The supplied original passage asks about the remaining Lie-type families and their shortest mixed and non-singular identities; no precise page within the report is supplied.
\item Supplied literature reference: Bradford--Schneider--Thom, finite mixed identities. \url{https://arxiv.org/abs/2306.14532}.
\item Supplied literature reference: Bradford--Ersoy--Schneider--Thom, 2026. \url{https://arxiv.org/abs/2608.14537}.
\item Supplied literature reference labeled ``explicit open questions'': \url{https://doi.org/10.4171/owr/2025/16}.
\item Literature assessment dated exactly \texttt{2026-08-22}: the supplied curation labels the problem partially solved and reports substantial bounds for the remaining Lie-type families and shortest non-singular identities, together with new 2026 structural results, while stating that the problem is not fully resolved for all families. This is supplied dated curation derived from the OWR collection, not independently verified source inspection or a new literature assessment. The record itself requests independent verification of time-sensitive status and cited claims.
\item The supplied statement carries the registry label \texttt{corrected\_verified}; its accompanying qualification says that only clear formula delimiters, notation, grammar, and typography were repaired and that mathematical content and scope were preserved. This label does not establish mathematical correctness or resolve the missing context.
\end{itemize}
\end{document}
